\begin{textsample}{POS Dim 1 – al – Score 48.00 – al/al\_em\_39.txt}  \label{ex:f1_pos_153}
SABERES ESCOLARES DOS COMPONENTES CURRICULARES DAS \textbf{CIÊNCIAS} HUMANAS E SOCIAIS APLICADAS Na primeira parte do texto \textbf{introdutório} expomos alguns dos principais pontos a respeito dos lugares semânticos, sociais, culturais e pedagógicos que as Ciências Humanas e Sociais Aplicadas ocupam na formação integral dos/as estudantes alagoanos/as.
A partir de agora vamos expor, de forma genérica, os principais saberes escolares dos componentes curriculares que compõem a nossa área de conhecimento.
Entendemos ser importante trazer aqui, como \textbf{salientamos} antes, uma discussão mínima sobre \textbf{métodos}, \textbf{abordagens} e objetos que fazem parte do fazer científico da Filosofia, da Geografia, da História e da Sociologia.
Mesmo sabendo que o diálogo interdisciplinar é, não somente \textbf{oportuno}, mas \textbf{constitutivo} do trabalho de ensino-aprendizado do Novo Ensino Médio, pensá -las na sua individualidade pode trazer um melhor auxílio e clareza aos docentes quanto ao entendimento do que se propõe neste documento.
Importa esclarecer que a \textbf{exposição} a seguir foca nas contribuições científicas de cada campo do saber para a constituição das Ciências Humanas Sociais Aplicadas, o que não significa que a prática docente deve ignorar os diversos outros saberes, sobretudo tomando -os como ponto de partida para a apresentação, de forma \textbf{dialógica}, dos conhecimentos científicos ainda não conhecidos pelos/as estudantes do ensino \textbf{médio}. \textbf{FILOSOFIA} O ensino de Filosofia no Brasil vem de uma longa trajetória.
Desde a chegada dos Jesuítas às terras do pau-brasil, tem sua história registrada em entradas e saídas do currículo escolar na educação brasileira.
Sua permanência ou não, variava de acordo com a política educacional e as correntes filosóficas vigentes ; e também as formas e relações de poder estabelecidas pelos governantes de cada época.
Na colonização, considerada fundamental para estabelecer o projeto católico contra reformista, na virada do \textbf{século} XIX para o XX, era vista como uma forma de estabelecer a racionalidade positivista.
Ao longo da primeira metade do \textbf{século} XX esteve presente no ensino secundário, sendo excluída no período da ditadura militar e retornando somente em 2008, quando da sua inclusão obrigatória no currículo juntamente com a Sociologia.
O retorno da Filosofia foi realizado após muitas lutas reivindicatórias dos professores, professoras e entidades de classe, e seu ensino escolar estimulou muitos \textbf{debates} dentro das Universidades sobre a democratização desse saber.
Seguindo o \textbf{ideal} socrático de que todos já somos seres filosóficos, o ensino de Filosofia faz -se necessário para que as reflexões e questões filosóficas possam vir à tona nos/as estudantes do ensino \textbf{médio}.
Construir o conhecimento, do ponto de \textbf{vista} filosófico ( e também etimológico ), tem a ver com a elaboração de um pensamento autônomo, de deixar ir além.
O filosofar, portanto, seria uma atividade de ensinar/aprender em si - por conta dessa perspectiva de conhecer os conceitos mais profundamente.
A presença do conhecimento filosófico na escola pretende, nesse sentido, ensinar a filosofar, buscando a construção e reconstrução de conceitos.
Não é trazer soluções e respostas, mas pensar a existência e a experiência, realizando o estímulo ao pensamento \textbf{crítico} e autônomo.
Dessa forma, a Filosofia tem o papel de dar complexidade ao óbvio, ao que o/a estudante conhece empiricamente, a priori.
Ela deixa de ser um conhecimento apenas abstrato e passa a ser vivenciada, experimentada em sua vida social, \textbf{econômica}, profissional, emocional, psicológica e etc.
A Filosofia pode mostrar aos/às estudantes o sentido de sua existência concreta, torna -se formativa ao ajuda -los a se dar conta do lugar que ocupa na realidade histórica objetiva, como ele ( ela ) se situa no contexto real de existência.
O papel pedagógico da Filosofia é \textbf{subsidiar} os/as estudantes a ler seu mundo, através da mediação curricular.
Ou seja, a Filosofia pode, através do exercício do pensar, buscar incentivar o/a estudante a entender seu lugar no mundo.
Um dos \textbf{métodos} mais frutíferos de ensinar a filosofar é a partir da pergunta filosófica.
Esta pressupõe, de acordo com Silveira ( 2011 ) uma espécie de “ sei que nada sei ” de Sócrates, como ponto de partida para a descoberta do saber.
O professor e a professora devem ser os mediadores com os/as estudantes para que eles possam dialogar com essa perspectiva de reconhecer a própria ignorância, buscando superá -la e \textbf{aprimorar} seu jeito de pensar.
Um segundo passo importante, é estimular os/as estudantes a adquirirem uma atitude problematizadora por meio de reflexões realizadas a partir de problemáticas vivenciadas pelos/as estudantes em sua realidade objetiva e que sejam \textbf{pertinentes} às diversas áreas da Filosofia.
Como exemplo, podemos citar questionamentos acerca da participação política ( se realmente acontece ), o significado efetivo da cidadania, desigualdade de classes ( o que são classes sociais?
Por que existem? ) como pontos a serem questionados e que se \textbf{aproximam} da Filosofia Política.
Também podemos falar de Ética ao problematizar os valores morais, as noções de Bem e Mal ; de Estética ao questionar o que é Belo ou Feio etc.
O professor precisa mediar o conhecimento, ao fazer com que o/a estudante perceba a necessidade de se refletir acerca dessas \textbf{problematizações}, mostrando que tudo isso faz parte do seu cotidiano, que são \textbf{problemas} que também lhes \textbf{dizem} respeito.
Nesse sentido, a partir do conhecimento filosófico é possível realizar associações com a realidade social do educando, ao buscar relacionar fenômenos que fazem parte do cotidiano à conceitos próprios da Filosofia.
Ao tratar da participação política, por exemplo, é possível resgatar os filósofos gregos como Platão e Aristóteles para que se possa compreender os primórdios da participação política e também as bases do que podemos compreender hoje como democracia.
Pode -se também partir da realidade social concreta e associá -la a questões filosóficas mais amplas e \textbf{conceituais}.
Para tratar do racismo, por exemplo, podemos pegar notícias recentes de jovens negros sendo assassinados como motivador para se debater questões filosóficas \textbf{relevantes} como alteridade, ética, violência, dentre outras.
Outro ponto \textbf{relevante} ao pensarmos as \textbf{aulas} de Filosofia, é o estímulo à leitura dos textos filosóficos, pois o desenvolvimento da \textbf{capacidade} de ler e entender tais textos nos leva a uma educação para o exercício da dúvida, da contestação e do pensamento.
A partir da leitura de Filósofos, podemos também estimular nossos/as estudantes a pensar a criação de conceitos, como nos incita Lima ( 2010 ), ao afirmar que a Filosofia é uma \textbf{disciplina} criadora de conceitos para que possamos nos situar melhor no espaço-tempo.
Ao debater temas fundamentais – presentes nas Ciências Humanas como um todo - como questões étnico-raciais, diversidade de gênero, violência, política, relações sociais, cidadania, liberdade, relações de poder, dentre outras \textbf{tantas} que estão presentes no cotidiano dos/as estudantes, podemos incentivá -los à construção de um pensamento \textbf{crítico} e autônomo, bem como a se reconhecerem como seres sociais e políticos no contexto e território em que estão inseridos.
Importa, nesse processo, chamar atenção para o \textbf{foco} na valorização das identidades quilombolas e indígenas presentes em Alagoas, com base na lei 11.645/08, que determina o ensino de história e cultura afro-brasileira e indígena.
GEOGRAFIA A Geografia é uma \textbf{disciplina} que estuda a relação do homem com o meio, isto é, com o espaço geográfico, nos \textbf{auxiliando} no entendimento do espaço vivido, formado por elementos históricos, sociais e naturais.
Ela é fundamental para a formação integral dos nossos/as estudantes, visto que possibilita pensar as transformações da natureza nas suas relações com as dinâmicas sociais.
De acordo com a OCEM ( 2006:44 ) a Geografia no Ensino Médio abre \textbf{caminhos} à múltiplas possibilidades de ensino e aprendizagem próprios dessa \textbf{ciência}.
Esse percurso favorece a percepção do espaço geográfico, e do papel ocupado pelo próprio/a estudante em uma sociedade que está constante transformação, favorecendo o desenvolvimento da leitura geográfica do mundo que o rodeia.
Além de orientar a formação cidadã no sentido de aprender a conhecer, aprender a fazer, aprender a conviver e aprender a ser, reconhecendo as contradições e os conflitos existentes no mundo.
Para OCEM ( 2006 : 44 ) um dos objetivos da Geografia no ensino \textbf{médio} é a organização de conteúdos que permitam ao/a estudante e a \textbf{aluna} experienciarem aprendizagens significativas através da percepção, apropriação, \textbf{análise}, discussão e proposição de \textbf{problemas} ambientais, sociais, da ocupação do espaço e das transições populacionais, fluxos migratórios, desigualdades \textbf{econômicas}, sociais e culturais.
Essa é uma concepção contida em \textbf{teorias} de aprendizagem que enfatizam a necessidade de considerar os conhecimentos prévios dos/as estudantes e o meio geográfico no qual estão inseridos, partindo destes para trabalhar os conceitos específicos como : lugar, espaço, tempo, região, território, natureza, paisagem e sociedade.
Dessa forma, a ação humana apresenta -se como elemento determinante das mudanças constantes do espaço, sociedade e natureza devem ser problematizadas dialeticamente.
De acordo com os PCNs ( 2000 : 109 ) : “ A observação, \textbf{descrição}, experimentação, analogia e síntese devem ser ensinadas para que os/as estudantes [ e as \textbf{alunas} ] possam aprender a explicar, compreender e até mesmo representar os processos de construção do espaço e dos diferentes tipos de paisagens e territórios ”.
Tal como as outras \textbf{ciências}, a Geografia passa por transformações constantes que se mostram como um desafio à prática docente.
Desta feita, a utilização de variadas fontes, o exercício constante e contextualizado da leitura dos espaços geográficos, consagrando o \textbf{debate} entre os/as estudantes como forma de \textbf{posicionamento} político são alguns apontamentos que devem ser levadas em consideração.
A Geografia deve permitir e possibilitar aos educandos uma \textbf{compreensão} da realidade que leve em conta a premissa de que o espaço geográfico, em todas as suas dimensões, se constrói em uma relação dinâmica com o espaço vivido.
Nessa perspectiva de ensino e aprendizagem da Geografia deve ter como premissas \textbf{metodológicas} a \textbf{interdisciplinaridade}, a \textbf{transdisciplinaridade} e a \textbf{contextualização}.
Por exemplo, é possível trabalhar fenômenos que inclui o homem como parte de toda a mudança e transformação do espaço geográfico utilizando -se escritores alagoanos como Graciliano Ramos com o livro “ Vidas Secas ” e Jorge de Lima com “ Calunga ” ou até mesmo a Literatura de Cordel, poemas, canções, produções de \textbf{vídeos} que abordem o espaço vivido pelo/as estudante.
Por meio destas obras, dentre outras possibilidades, podem -se conhecer os diferentes aspectos regionais, sociais, culturais e naturais de uma região.
Pode -se também, ao analisar a hidrosfera e a biosfera utilizar -se de algum embasamento dos componentes curriculares de Química, Matemática e Biologia, como por exemplo o estudo de tabelas, gráficos e da linguagem \textbf{cartográfica}, bem como trazer para o \textbf{debate} os saberes locais sobre esses espaços.
A BNCC apresenta a ideia de que o/a estudante deve ser o protagonista deste espaço geográfico. \textbf{Dito} isto, apontamos que essa percepção da Geografia deve contribuir para o desenvolvimento de várias \textbf{competências} \textbf{socioemocionais} como empatia, resolução de \textbf{problemas}, comunicação, visão \textbf{crítica} da sociedade dentre outras favorecendo enfim a formação integral e cidadã.

% matched lemmas: abordagem, aluno, análise, aprimorar, aproximar, aula, auxiliar, caminho, capacidade, cartográfico, ciência, competência, compreensão, conceitual, constitutivo, contextualização, crítico, debate, descrição, dialógico, disciplina, dizer, econômico, exposição, filosofia, foco, ideal, interdisciplinaridade, introdutório, metodológico, médio, método, oportuno, pertinente, posicionamento, problema, problematização, relevante, salientar, socioemocional, subsidiar, século, tanto, teoria, transdisciplinaridade, vista, vídeo
\end{textsample}
